\title{FlashAttention-3: Fast and Accurate Attention with Asynchrony and Low-precision}

\begin{abstract}
Within the prevalent Transformer framework, attention serves as a fundamental mechanism but remains a key performance constraint in large language models and tasks demanding extensive context windows. \textsc{FlashAttention} previously introduced a strategy to accelerate attention on GPUs by optimizing memory access patterns. Yet, this approach does not fully exploit recent advancements in hardware capabilities; notably, \textsc{FlashAttention-2} reaches only 35\% computational utilization on the H100 GPU. In this work, we introduce three principal innovations to enhance attention performance on Hopper GPUs: leveraging the concurrent execution features of Tensor Cores and TMA to (1) overlap computation with data transfer using warp-specialization, (2) concurrently perform block-level matrix multiplication and softmax operations, and (3) utilize block quantization alongside incoherent processing, capitalizing on native FP8 low-precision support. Through these methods, our system, \textsc{FlashAttention-3}, demonstrates a 1.5--2.0$\times$ acceleration on H100 GPUs, with FP16 computations attaining up to 740 TFLOPs/s (75\% utilization), and FP8 peaking near 1.2 PFLOPs/s. Moreover, empirical results show that FP8 \textsc{FlashAttention-3} achieves 2.6$\times$ less numerical error relative to a standard FP8 attention implementation.
\end{abstract}

\section{Introduction}
\label{sec:intro}

Within the Transformer architecture~\citep{vaswani2017attention}, the main computational challenge arises from the attention mechanism, particularly due to the quadratic time complexity of calculating self-attention scores between queries and keys as the sequence length increases. Enabling attention mechanisms to efficiently handle much longer contexts stands to provide substantial benefits, such as improved modeling and reasoning capabilities across multiple extended documents~\citep{guo2021longt5,shaham2022scrolls,peng2023yarn} and large-scale codebases~\citep{roziere2023code, li2023starcoder}; the ability to process new data modalities like detailed images~\citep{chen2022scaling}, audio~\citep{gulati2020conformer}, or video~\citep{ho2022video}; as well as expanding applications including long-range user interactions~\citep{sun2019bert4rec} and agents engaged in complex workflows over extended time frames~\citep{yao2022react}. Consequently, considerable research attention has focused on accelerating attention computation for long contexts, utilizing a variety of strategies such as approximation algorithms~\citep{katharopoulos2020transformers,choromanski2020rethinking,
tay2020efficient}, software-level enhancements (\citep{rabe2021self,dao2022flashattention,kwon2023efficient}), and even the development of novel alternative architectures~\citep{peng2023rwkv,sun2023retentive,gu2023mamba}.

This study extends the approach introduced by \citet{dao2022flashattention}, which emphasizes designing exact-attention algorithms by aligning them with the GPU's execution mechanisms and hardware properties. The original \textsc{FlashAttention} method presented in \citep{dao2022flashattention} by Dao et al. leverages an innovative tiling methodology to enhance parallelism in attention computation. By consolidating all attention-related operations within a unified GPU kernel, the approach avoids unnecessary accesses to slow global memory during intermediate stages.

Subsequently, \citet{dao2023flashattention2} redesigned this algorithm as \textsc{FlashAttention-2}, incorporating parallelization across the sequence length axis and structuring the forward pass's innermost loop to operate over blocks of the key and value matrices. These changes serve to improve GPU resource occupancy and distribute computational tasks more efficiently. Nevertheless, our evaluation indicates that \textsc{FlashAttention-2} continues to display suboptimal GPU utilization—especially on more advanced hardware, such as the Hopper H100, where usage rates lag behind high-performance matrix-multiplication (GEMM) kernels (35\% versus 80–90\% utilization). Some of this discrepancy is attributable to implementation nuances, notably the continued reliance on Ampere-era instructions instead of exploiting Hopper-specific features on Tensor Cores. Recent advances, such as ThunkerKitten~\citep{spector2024thunder} and cuDNN 9~\citep{cudnn9}, demonstrate that applying Hopper-targeted instructions alongside tile-oriented abstractions not only accelerates attention operations but also streamlines their coding and maintenance.

At its core, the algorithm underpinning \textsc{FlashAttention-2} operates within a straightforward synchronous paradigm and does not deliberately integrate asynchronous computation or reduced numerical precision in its architecture. Asynchronous execution arises from hardware advancements designed to boost key machine learning tasks: specialized components are dedicated to matrix multiplications (Tensor Cores) or memory transfers (Tensor Memory Accelerator, abbreviated as TMA), while other CUDA cores focus on tasks involving logic, integer arithmetic, or floating-point calculations. The adoption of low-precision formats—such as FP8 in Hopper, FP4 in Blackwell, following earlier developments like FP16 (introduced on Pascal in 2017) and BF16 (featured in Ampere in 2020)—enables substantially increased computational throughput (2x or 4x) without greater energy or silicon requirements. The innovations in Hopper relevant to these topics are examined in \cref{subsec:hardware}. 

The key technical obstacle lies in reengineering \textsc{FlashAttention-2} so that it leverages these hardware capabilities: supporting asynchrony necessitates concurrent execution of matmul and softmax computations, despite their data dependencies; adopting low-precision approaches demands careful handling to curtail quantization error, which becomes particularly important for feature outliers within LLMs~\citep{dettmers2208llm, sun2024massive}.

For this purpose, we introduce \textsc{FlashAttention-3}, incorporating three innovative concepts that collectively advance efficiency on modern GPU platforms:\footnote{Our experiments focus on NVIDIA's Hopper architecture, but the proposed algorithm applies broadly to any GPU system offering strong support for asynchronous operations and low-precision computation.}

\begin{enumerate}

\item \textbf{Producer-Consumer asynchrony:} We introduce a warp-specialized software pipelining approach that leverages the decoupled execution of data transfers and Tensor Core computations. By assigning the roles of data producers and consumers to separate warps, this method enhances the algorithm’s capacity to conceal delays caused by memory access and instruction issuance.

\item \textbf{Hiding softmax under asynchronous block-wise GEMMs:} We concurrently execute the non-GEMM operations of softmax, such as floating point multiply-add and exponentiation—which are bandwidth-limited—with asynchronous WGMMA GEMM operations. To facilitate this overlap, we revise the \textsc{FlashAttention-2} algorithm to minimize sequential constraints between softmax and GEMMs. For instance, in our algorithm’s two-stage structure, one block of the score matrix undergoes softmax processing while, at the same time, WGMMA operates in an asynchronous manner to compute the subsequent block.

\item \textbf{Hardware-accelerated low-precision GEMM:} The forward pass algorithm is modified to enable the use of FP8 Tensor Cores for GEMM, resulting in a near doubling of measured TFLOPs/s. This adaptation entails resolving the distinct memory layout requirements for WGMMA, specifically with respect to how blocks of FP32 accumulators and FP8 operand matrices are structured in memory. Block quantization and incoherent processing techniques are employed to reduce the loss of numerical precision associated with FP8 computation.

To empirically substantiate our approach, we conduct benchmarking of \textsc{FlashAttention-3} on the H100 SXM5 GPU across varying parameter settings. Our findings indicate the following: (1) with FP16, the forward computation experiences a speedup ranging from 1.5 to 2.0$\times$ relative to \textsc{FlashAttention-2}, attaining up to 740 TFLOPs/s, while the backward computation benefits from a 1.5-1.75$\times$ improvement; (2) employing FP8, the performance reaches nearly 1.2 PFLOPs/s; and (3) for large sequence lengths, FP16 delivers superior results and FP8 demonstrates competitiveness\footnote{Specifically, for a head dimension of 64, \textsc{FlashAttention-3} FP8 surpasses others, whereas for head dimensions of 128 and 256 it matches performance in the absence of causal masking, but lags when causal masking is present.} when compared against NVIDIA's cuDNN library’s attention implementation. Additionally, we confirm that FP16 \textsc{FlashAttention-3} exhibits numerical errors on par with \textsc{FlashAttention-2} and provides improved accuracy over conventional attention approaches, as all intermediate values such as those from softmax rescaling are stored in FP32. In cases involving outlier features, FP8 \textsc{FlashAttention-3} configured with block quantization and incoherent processing achieves 2.6$\times$ greater accuracy than standard attention utilizing per-tensor quantization.

We have released \textsc{FlashAttention-3} as open-source under a flexible license\footnote{\textsc{FlashAttention-3}
  can be found at \texttt{https://github.com/Dao-AILab/flash-attention}} and intend to incorporate it into both PyTorch and Hugging Face libraries, aiming to maximize its accessibility and utility for the research and development communities.

\section{Background: Multi-Head Attention and GPU Characteristics}
\label{sec:background}

\subsection{Multi-Head Attention}
\label{subsec:multi_head_attn}

Let $\mathbf{Q}, \mathbf{K}, \mathbf{V} \in \mathbb{R}^{N \times d}$ be the query, key and value input sequences associated to a single head, where $N$ is the sequence length and $d$ is the head dimension. Then the attention output $\mathbf{O}$ is computed as:

\begin{equation*}
  \mathbf{S} = \alpha \mathbf{Q} \mathbf{K}^\top \in \mathbb{R}^{N \times N}, \quad \mathbf{P} = \mathrm{softmax}(\mathbf{S}) \in \mathbb{R}^{N \times N}, \quad \mathbf{O} = \mathbf{P}\mathbf{V} \in \mathbb{R}^{N \times d},
\end{equation*}

The $\mathrm{softmax}$ operation is performed on each row, and it is common to choose $\alpha = 1/\sqrt{d}$ as the scale parameter. To avoid issues with numerical stability arising from exponentiation, we usually subtract $\mathrm{rowmax}(\mathbf{S})$ from $\mathbf{S}$ before applying further calculations. In the case of multi-head attention (MHA), distinct query, key, and value projection matrices are used for each attention head. These operations are executed concurrently across several heads and batches, yielding the aggregated output tensor.

Now let $\phi$ be a scalar loss function and let $\mathbf{d}(-) = \partial \phi / \partial (-)$ be notation for the gradient. Given the output gradient $\mathbf{dO} \in \mathbb{R}^{N \times d}$, we compute $\mathbf{dQ}$, $\mathbf{dK}$, and $\mathbf{dV}$ according to the chain rule as follows:

\begin{align*}
  \mathbf{dV} &= \mathbf{P}^\top \mathbf{dO} \in \mathbb{R}^{N \times d} \\
  \mathbf{dP} &= \mathbf{dO} \mathbf{V}^\top \in \mathbb{R}^{N \times N} \\
  \mathbf{dS} &= \mathrm{dsoftmax} (\mathbf{dP}) \in \mathbb{R}^{N \times N} \\
  \mathbf{dQ} &= \alpha \mathbf{dS} \mathbf{K} \in \mathbb{R}^{N \times d} \\
  \mathbf{dK} &= \alpha \mathbf{dS}^\top \mathbf{Q} \in \mathbb{R}^{N \times d},
\end{align*}

In this context, $\mathbf{d}s = (\mathrm{diag}(p) - p p^\top)\mathbf{d}p$ holds where $p = \mathrm{softmax}(s)$ and $s$ represents a vector input. We denote $\mathrm{dsoftmax}(\mathbf{dP})$ as this operation performed on each row. As before, this calculation can be executed in parallel over multiple attention heads and batches during the backward phase of MHA.

\subsection{GPU hardware characteristics and execution model}
\label{subsec:hardware}

We discuss the features of the GPU execution model pertinent to \textsc{FlashAttention-3}, concentrating specifically on how these features are implemented in the NVIDIA Hopper architecture.

\paragraph{Memory hierarchy:} On GPUs, the memory system is structured as a layered hierarchy, where higher storage capacities are coupled with lower bandwidths (\cref{tab:gpu-hierarchy})\footnote{\citet{luo2024benchmarking} provides a shared memory bandwidth of 128 bytes per clock cycle per SM, which we scale by 132 SMs and a boost clock speed of 1830 MHz.}. The largest of these layers, global memory (GMEM) or HBM, resides off-chip and can be accessed by all streaming multiprocessors (SMs). GMEM data is automatically cached within an on-chip L2 cache to improve locality. Further within the hierarchy, each SM is equipped with shared memory (SMEM)—a small, on-chip cache with high bank parallelism that is explicitly managed by programmers. At the base of the hierarchy, each SM has a local register file for the fastest data access.

\paragraph{Thread hierarchy:} The GPU utilizes a structured programming model based on hierarchical assemblies of threads. Beginning with individual threads at the lowest level, the hierarchy consists of threads, followed by warps (collections of 32 threads), warpgroups (each made up of 4 consecutive warps), then threadblocks (also referred to as cooperative thread arrays or CTAs), threadblock clusters (specific to Hopper architectures), and finally grids at the highest level.

There is a strong interdependence between these two hierarchical structures. All threads belonging to a single CTA are executed concurrently on the same SM, and multiple CTAs assigned to one cluster are executed on the same GPC. The SMEM can be accessed by every thread within a CTA, while each individual thread possesses its own set of up to 256 private registers (RMEM).

\paragraph{Asynchrony and warp-specialization:}

GPUs are designed as throughput-oriented processors that leverage high levels of parallelism and asynchronous operations to mask the latencies associated with memory access and execution. In the context of asynchronous memory transfers between GMEM and SMEM, Hopper introduces the Tensor Memory Accelerator (TMA), a specialized hardware module \cite[\S7.29]{cuda}. Additionally, diverging from earlier architectures like Ampere, Hopper's Tensor Core—accessed through the warpgroup-wide WGMMA instruction \cite[\S9.7.14]{ptx}—operates asynchronously and is capable of reading inputs straight from shared memory.

Asynchronous hardware support enables the creation of warp-specialized kernels, wherein individual warps within a CTA are assigned distinct roles—either as producers, focusing solely on data movement, or as consumers, dedicated exclusively to computation. This division of labor enhances the compiler's capability to construct more effective instruction schedules \citep{warp-specialization-2011}. Moreover, the Hopper architecture facilitates the dynamic adjustment of register allocation among warpgroups using the \verb|setmaxnreg| instruction \cite[\S9.7.17.1]{ptx}, permitting warps performing MMAs to be granted a larger portion of RMEM compared to those restricted to TMA operations (which only require one thread).

\paragraph{Low-precision number formats:}
\label{sec:low-precision-gpu}
Contemporary GPUs incorporate dedicated hardware modules designed to expedite calculations using low-precision formats. As an illustration, the WGMMA instruction leverages FP8 Tensor Cores on Hopper architecture, resulting in a performance boost that achieves twice the throughput per SM in comparison to FP16 or BF16.

To use FP8 WGMMA effectively, one must be aware of the operand layout requirements it imposes. For a GEMM operation computing $A \times B^{\top}$ where $A$ is an $M\times K$ matrix and $B$ an $N\times K$ matrix, we characterize each operand as \emph{mn-major} if elements are stored contiguously along the $M$ or $N$ dimension, and as \emph{k-major} if contiguity is along the $K$ dimension. FP16 WGMMA offers flexibility, permitting operands in SMEM to be either mn-major or k-major. In contrast, FP8 WGMMA restricts input in SMEM to the k-major layout only. Additionally, in scenarios such as attention mechanisms where sequential GEMM computations are fused into a single kernel, mismatches between the FP32 accumulator layout and FP8 operand arrangement complicate the chaining of FP8 WGMMAs, making it challenging to execute dependent operations efficiently.

Within the framework of attention mechanisms, these layout constraints necessitate specific adjustments to the architecture of an FP8 algorithm. The details of these modifications are provided in \cref{sec:algofp8}.

\subsection{Standard Attention and Flash Attention} In line with \citet{dao2022flashattention}, we refer to \textbf{standard attention} as the GPU implementation of attention that writes the intermediate matrices $\mathbf{S}$ and $\mathbf{P}$ to high-bandwidth memory (HBM). The primary innovation behind \textsc{FlashAttention} is to utilize a local softmax reduction, which allows it to circumvent costly memory accesses for these intermediates and instead combine all attention computations within a single kernel operation. The local softmax operation is realized by lines \ref{code-ws:softmax_start}-\ref{code-ws:softmax_end} of the consumer mainloop presented in \cref{alg:flash3_wgmma_ws_only}, alongside the normalization steps applied to segments of $\mathbf{O}$. A straightforward proof showing that this approach yields the correct computation of $\mathbf{O}$ is available in \cite[\S 2.3.1]{dao2023flashattention2}.

\section{FlashAttention-3: Algorithm}
\label{sec:algo}

This section introduces the \textsc{FlashAttention-3} algorithm, concentrating on the forward computation; details regarding the backward computation can be found in~\cref{sec:algo_ws_bwd}. Initially, we demonstrate the method for incorporating warp-specialization and a circular SMEM buffer into the foundational \textsc{FlashAttention-2} algorithm. Next, we outline the use of WGMMA asynchrony to construct a two-stage pipeline that overlaps GEMM and softmax operations. Lastly, we present the necessary adjustments for FP8, addressing both layout compatibility and accuracy improvements through block quantization and incoherent processing.

\subsection{Producer-Consumer asynchrony through warp-specialization and pingpong scheduling}
\label{sec:algo_ws}

\paragraph{Warp-specialization}
Similar to the approach in \textsc{FlashAttention-2}, the forward computation in \textsc{FlashAttention-3} naturally lends itself to parallelization across the batch dimension, head indices, and the query sequence axis. For this reason, an overview at the CTA level suffices: here, a tile $\mathbf{Q}_i$ from the query matrix is processed to yield the matching segment $\mathbf{O}_i$ of the output. To streamline the exposition, we describe the warp-specialized design utilizing a circular SMEM buffer, but excluding the added complexity of GEMM-softmax fusion. Denote the head size by $d$, sequence length as $N$, and select a query block size $B_r$. The query matrix $\mathbf{Q}$ is then partitioned into $T_r = \lceil \frac{N}{B_r} \rceil$ blocks, labeled as $\mathbf{Q}_1, .., \mathbf{Q}_{T_r}$.

\begin{algorithm}[H]
    \caption{\small\label{alg:flash3_wgmma_ws_only}\textsc{FlashAttention-3} forward pass \textbf{without} intra-consumer overlapping -- CTA perspective}
    \begin{algorithmic}[1]
\REQUIRE Matrices $\mathbf{Q}_i \in \mathbb{R}^{B_r \times d}$, $\mathbf{K}, \mathbf{V} \in \mathbb{R}^{N \times d}$ residing in HBM, key block size $B_c$, and compute $T_c = \lceil \frac{N}{B_c} \rceil$.
\STATE Set up a pipeline controller to coordinate barrier synchronization, utilizing an $s$-stage circular SMEM buffer.
\IF {executing within the producer warpgroup}
\STATE Release a specified set of registers in advance.
\STATE Trigger load of $\mathbf{Q}_i$ from HBM into shared memory.
\STATE When this load completes, signal consumers regarding the availability of $\mathbf{Q}_i$.
\FOR{$0 \le j < T_c$}
    \STATE Await confirmation that the $(j\,\%\,s)$th buffer stage has been processed and freed.
    \STATE Load $\mathbf{K}_j$ and $\mathbf{V}_j$ from HBM into the appropriate buffer stage $(j\,\%\,s)$ in shared memory.
    \STATE Once the load finishes, notify consumers that $\mathbf{K}_j$ and $\mathbf{V}_j$ are ready.
\ENDFOR
\ELSE \STATE Claim a reserved number of registers, adjusted by the count of consumer warps.
\STATE Locally, allocate $\mathbf{O}_i = (0) \in \mathbb{R}^{B_r \times d}$, $\ell_i, m_i = (0), (-\infty) \in \mathbb{R}^{B_r}$.
\STATE Wait for $\mathbf{Q}_i$'s arrival in shared memory before proceeding.
\FOR{$0 \le j < T_c$}
\STATE Wait for $\mathbf{K}_j$ to be staged in shared memory.
\STATE Calculate $\mathbf{S}_i^{(j)} = \mathbf{Q}_i \mathbf{K}_j^T$ via SS-GEMM. Then commit and synchronize. \label{alg:ws_only_gemm1}
\STATE Preserve $m_i^{\mathrm{old}} = m_i$; update $m_i = \mathrm{max}(m_i^{\mathrm{old}}, \mathrm{rowmax}(\mathbf{S}_i^{(j)}))$. \label{code-ws:softmax_start}
\STATE Update $\widetilde{\mathbf{P}}_i^{(j)} = \mathrm{exp}(\mathbf{S}_i^{(j)} - m_i)$ and adjust $\ell_i = \mathrm{exp}(m_i^{\mathrm{old}} - m_i) \ell_i + \mathrm{rowsum}(\widetilde{\mathbf{P}}_i^{(j)})$. \label{code-ws:softmax_end}
\STATE Hold for $\mathbf{V}_j$ to become available in shared memory.
\STATE Perform $\mathbf{O}_i = \mathrm{diag}(\mathrm{exp}(m_i^{\mathrm{old}} - m_i))^{-1} \mathbf{O}_i + \widetilde{\mathbf{P}}_i^{(j)} \mathbf{V}_j$ through RS-GEMM. Commit and synchronize. \label{alg:ws_only_gemm2}
\STATE Mark the $(j\,\%\,s)$th buffer stage as complete, permitting the producer to proceed.
\ENDFOR
\STATE Finalize by calculating $\mathbf{O}_i = \mathrm{diag}(\ell_i)^{-1} \mathbf{O}_i$ and $L_i = m_i + \log(\ell_i)$.
\STATE Output $\mathbf{O}_i$ and $L_i$ back to HBM to represent the $i$th segment of $\mathbf{O}$ and $L$.
\ENDIF
\end{algorithmic}
\end{algorithm}

In our deployment of \cref{alg:flash3_wgmma_ws_only} on Hopper, we leverage \verb|setmaxnreg| for managing (de)allocations, employ TMA for retrieving $\mathbf{Q}_i$ as well as $\{ \mathbf{K}_j, \mathbf{V}_j \}_{0 \leq j < T_c}$, and utilize WGMMA to perform GEMM operations within the consumer mainloop. The SS and RS prefixes distinguish whether the initial operand is loaded from shared memory or the register file, respectively. When analyzing the control flow of \cref{alg:flash3_wgmma_ws_only}, it is important to recognize that TMA load instructions operate asynchronously and thus are not blocked by ongoing loads. Additionally, within the producer mainloop, for the first $s$ iterations, no waits are triggered since the buffer is in the process of being populated.

\paragraph{Pingpong scheduling} Because WGMMA and TMA operate asynchronously and use warp-specialization, it becomes feasible to concurrently execute the softmax computation in one warpgroup while another warpgroup carries out the GEMM. This is particularly advantageous given that, on current hardware accelerators, matmul operations are orders of magnitude more performant than non-matmul operations. For instance, the H100 SXM5 GPU achieves 989 TFLOPS for FP16 matmul but provides only 3.9 TFLOPS for special functions like exponential\footnote{According to the CUDA programming guide, each streaming multiprocessor (SM) is able to perform 16 special function operations per clock cycle. With 132 SMs operating at 1830 MHz, this amounts to 3.9 TFLOPS for special function throughput.}, which are required for softmax. Examining the FP16 attention forward pass with a head dimension of 128, we see that matmul requires 512 times more FLOPS relative to exponential operations. Despite this, the throughput for exponential is 256 times less than that for matmul, which means the exponential step can consume as much as 50\% of the total cycle time. The disparity increases with FP8 precision, where matmul throughput doubles without any corresponding increase in exponential throughput.

The exponential computation is executed by a dedicated hardware component—the multi-function unit. Ideally, we aim to overlap this exponential operation with the matrix multiplication handled by the Tensor Cores. To orchestrate this, synchronization barriers (\texttt{bar.sync} instructions) are employed, ensuring that the GEMMs (specifically, GEMM1 -- $\mathbf{P} \mathbf{V}$ for the current iteration and GEMM0 -- $\mathbf{Q} \mathbf{K}^\top$ for the subsequent iteration) of warpgroup 1 are dispatched prior to those in warpgroup 2. Consequently, warpgroup 1 can conduct its softmax computations while warpgroup 2 executes its GEMMs. The process alternates thereafter, with warpgroup 2 performing the softmax as warpgroup 1 runs its GEMMs, leading to a "pingpong" scheduling pattern, as depicted in~\cref{fig:pingpong_scheduling}. While in real workloads this pingponging does not achieve perfect overlap as shown in the illustration, it nonetheless tends to yield a significant performance benefit; for instance, for FP16 forward passes with a head dimension of 128 and sequence length of 8192, throughput can increase from 570 TFLOPS to the range of 620-640 TFLOPS.

\paragraph{Attention variants} In handling multi-query attention~\citep{shazeer2019fast} as well as grouped query attention~\citep{ainslie2023gqa}, our methodology aligns with \textsc{FlashAttention-2}, modifying the tensor indexing scheme to prevent multiple copies of $\mathbf{K}$ and $\mathbf{V}$ from being stored in HBM.

\subsection{Intra-warpgroup overlapping GEMMs and softmax}

It is possible to achieve overlap between certain instructions from the softmax computation and those from the GEMMs, even when confined to a single warpgroup. Here, we outline a specific method to accomplish this.

In the attention algorithm, operations within the inner loop (main loop) have sequential dependencies that impede parallelization within a single iteration. For example, (local) softmax (lines \ref{code-ws:softmax_start} to \ref{code-ws:softmax_end}) relies on the output $\mathbf{S}_i^{(j)}$ of the first GEMM, while the second GEMM takes its result $\widetilde{\mathbf{P}}_i^{(j)}$ as an operand. Indeed, the wait statements in lines \ref{alg:ws_only_gemm1} and \ref{alg:ws_only_gemm2} of \cref{alg:flash3_wgmma_ws_only} serialize the execution of softmax and GEMMs. However, we can break these dependencies by pipelining across iterations through additional buffers in registers. Pursuing this idea, we propose the following two-stage\footnote{Note that the number of stages of the overlapping scheme is bounded by, but need not equal, the number $s$ of stages in the circular SMEM buffer.} GEMM-softmax pipelining algorithm:

\begin{algorithm}[H]
    \caption{\small\label{alg:flash3_wgmma}\textsc{FlashAttention-3} consumer warpgroup forward pass}
    \begin{algorithmic}[1]
      \REQUIRE  Matrices $\mathbf{Q}_i \in \mathbb{R}^{B_r \times d}$ and $\mathbf{K}, \mathbf{V} \in \mathbb{R}^{N \times d}$ stored in HBM; key block size $B_c$; $T_c = \lceil \frac{N}{B_c} \rceil$.
      \STATE Allocate a designated set of registers depending on the current count of consumer warps.
      \STATE Set up on-chip initializations for $\mathbf{O}_i = (0) \in \mathbb{R}^{B_r \times d}$ and $\ell_i, m_i = (0), (-\infty) \in \mathbb{R}^{B_r}$.
      \STATE Ensure that $\mathbf{Q}_i$ and $\mathbf{K}_0$ are fully present in shared memory before proceeding.
      \STATE Calculate $\mathbf{S}_{\mathrm{cur}} = \mathbf{Q}_i \mathbf{K}_0^T$ using WGMMA, and execute commit with synchronization.
      \STATE Open up the $0$th buffer stage associated with $\mathbf{K}$.
      \STATE Derive $m_{i}$, $\tilde{\mathbf{P}}_{\mathrm{cur}}$, and $\ell_{i}$ from $\mathbf{S}_{\mathrm{cur}}$, and adjust the scale for $\mathbf{O}_i$.
      \FOR{$1 \le j < T_c - 1$}
        \STATE Confirm $\mathbf{K}_j$ is accessible in shared memory.
        \label{code:mainloop_start}
        \STATE Calculate $\mathbf{S}_{\mathrm{next}} = \mathbf{Q}_i \mathbf{K}_{j}^T$ utilizing WGMMA; commit operation but do not block.
        \label{code:first_wgmma}
        \STATE Pause until $\mathbf{V}_{j-1}$ has been transferred into shared memory.
        \STATE Update $\mathbf{O}_{i} = \mathbf{O}_{i} + \tilde{\mathbf{P}}_{\mathrm{cur}} \mathbf{V}_{j-1}$ with WGMMA, commit but proceed without waiting.
        \label{code:second_wgmma}
        \STATE Synchronize after the WGMMA execution of $\mathbf{Q}_i \mathbf{K}_{j}^T$.
        \STATE Calculate $m_{i}$, $\tilde{\mathbf{P}}_{\mathrm{next}}$, and $\ell_{i}$ based on $\mathbf{S}_{\mathrm{next}}$. 
        \label{code:softmax}
        \STATE Await completion of $\tilde{\mathbf{P}}_{\mathrm{cur}} \mathbf{V}_{j-1}$ through WGMMA, then rescale $\mathbf{O}_i$.
        \STATE Unlock the buffer stage for $\mathbf{K}$ at position $(j\,\%\,s)$ and $\mathbf{V}$ at $(j-1\,\%\,s)$.
        \STATE Transfer the contents of $\mathbf{S}_{\mathrm{next}}$ into $\mathbf{S}_{\mathrm{cur}}$.
        \label{code:mainloop_end}
      \ENDFOR \STATE Wait for $\mathbf{V}_{T_c - 1}$ to become available in shared memory.
      \STATE Apply
      $\mathbf{O}_{i} = \mathbf{O}_{i} + \tilde{\mathbf{P}}_{\mathrm{last}} \mathbf{V}_{T_c - 1}$ using WGMMA and commit synchronously.

\STATE Epilogue: Adjust $\mathbf{O}_{i}$ proportionally according to $m_{i}$. Determine $L_{i}$ using both $m_{i}$ and $\ell_{i}$.
Store $\mathbf{O}_{i}$ and $L_{i}$ in HBM, placing them in the $i$-th segment of $\mathbf{O}$ and $L$.

\cref{alg:flash3_wgmma} substitutes for the consumer segment of \cref{alg:flash3_wgmma_ws_only}, thereby composing the full \textsc{FlashAttention-3} procedure tailored to FP16 computations. Conceptually, we employ the term WGMMA to denote asynchronous GEMM routines. Throughout the mainloop (see lines \ref{code:mainloop_start}–\ref{code:mainloop_end}), the execution of the second WGMMA in the $j$th iteration (line \ref{code:second_wgmma}) coincides with the softmax step associated with iteration $j+1$ (line \ref{code:softmax}), effectively overlapping the two operations.

Although the pipelined model shown above promises significant theoretical improvements, several pragmatic issues must be addressed:

\paragraph{Compiler reordering} The pseudocode outlines an optimal sequence of operations, but the NVCC compiler typically rearranges instructions to maximize efficiency. Such optimizations can interfere with the intended pipelining arrangement of WGMMA and non-WGMMA tasks, possibly resulting in unintended side effects or reduced performance benefits. Examination of the SASS output reveals that the compiler does produce overlapped instructions as anticipated (refer to Section~\ref{sec:2-stage-sass}).

\paragraph{Register pressure} Achieving high efficiency requires keeping register spilling to a minimum. The 2-stage pipeline design, however, introduces extra requirements: more registers are needed to store intermediate computations and preserve state between pipeline stages. Specifically, an additional $\mathbf{S}_{\mathrm{next}}$ value is stored in registers, increasing per threadblock register allocation by $B_r \times B_c \times \text{sizeof}(\text{float})$. This extra demand can limit the ability to use larger block sizes—an optimization strategy that also heavily consumes registers. Therefore, optimal choices must be guided by profiling and benchmarking.

\paragraph{3-stage pipelining} Building on the earlier 2-stage method, we develop a 3-stage pipeline variant designed to further overlap the second WGMMA operation with the softmax computation. This design has the potential to push Tensor Core utilization even higher, but comes at the cost of further increased register usage per threadblock, as the pipeline now contains an additional stage. Balancing the competing needs of tile size and pipeline length becomes more complex. Details and evaluations of this 3-stage pipeline appear in ~\cref{sec:3-stage}.

\subsection{Low-precision with FP8}
\label{sec:algofp8}

\textbf{Efficiency: layout transformations.} Executing the forward computation of \textsc{FlashAttention-3} using FP8 precision introduces distinct complications related to layout consistency that are not present when employing FP16.

To begin, we observe that the input tensors $\mathbf{Q}$, $\mathbf{K}$, and $\mathbf{V}$ are generally organized with the head dimension being contiguous. However, for the second GEMM step in FP8 WGMMA, enforcing the k-major requirement necessitates that $\mathbf{V}$ (specifically, the tiles of $\mathbf{V}$ placed into SMEM) are contiguous along the sequence length dimension. Since TMA load operations are unable to alter the contiguous axis, we are presented with two possible solutions: (1) perform a transpose on $\mathbf{V}$ in GMEM as a pre-processing action, or (2) execute a tile-wise transpose of $\mathbf{V}$ within the kernel, immediately after loading to SMEM. For approach (1), the transpose may be (1a) incorporated into the epilogue of an upstream operation such as the rotary embedding, or (1b) implemented via an isolated pre-processing transpose kernel\footnote{Highly optimized transpose kernels can attain speeds close to the device’s peak bandwidth~\citep{colfax_cutlass_transpose_2024}.}, exchanging the strides of sequence length and head dimensions. Nevertheless, (1a) presents challenges for integration with widely used libraries, and (1b) becomes inefficient in contexts where memory bandwidth is the primary limitation, such as inference.

For FP8 \textsc{FlashAttention-3}, we choose to implement strategy (2). The in-kernel transposition leverages the LDSM (\verb|ldmatrix|) and STSM (\verb|stmatrix|) instructions, which enable a group of threads within a warp to collaboratively move data from SMEM to RMEM and back at 128-byte blocks.\footnote{Although the PTX specification defines LDSM/STSM as operating on $8 \times 8$ arrays of 16-bit values \cite[\S 9.7.13.4.15-16]{ptx}, we can concatenate pairs of 8-bit elements to exploit these instructions in the FP8 format. That said, the transpose variants of LDSM/STSM do not divide packed 8-bit values, which introduces the need for specific register shuffles between LDSM and STSM operations to achieve a tile-level transpose; we skip further technical details here.} LDSM/STSM not only facilitate efficient register use—enabling computation in the producer warpgroup—but also support layout transposition during memory transfers. Additionally, after the initial round, we schedule the transpose operation for the upcoming $\mathbf{V}$ tile to overlap with the execution of the two WGMMAs operating on the previous $\mathbf{V}$ and the current $\mathbf{K}$ tile.

Secondly, we note that the FP32 accumulator used in an FP8 WGMMA does not share the same memory layout as operand A when it resides in registers, contrasting with FP16. \cref{fig:rmem_accum} and \cref{fig:rmem_operand} illustrate segments of these two layouts, showing the order in which entries are stored per thread in registers. By employing byte permute instructions, the accumulator from the first WGMMA can be reorganized into a structure compatible for input into a subsequent WGMMA, and aligning with the format of the $\mathbf{V}$ tile resulting from the in-kernel transpose. More precisely, referencing \cref{fig:rmem_accum}, we systematically rearrange the register order to
$$\{ \verb|d0 d1 d4 d5 d2 d3 d6 d7| \},$$
repeating this pattern across every 8 bytes. From the perspective of the logical layout of the $\mathbf{P}$ tile, this technique effectively permutes its columns; for instance, columns $0189$ are now positioned as the first four columns. To ensure WGMMA generates the correct output tile, the in-kernel transpose can be set to output a corresponding row permutation for the $\mathbf{V}$ tile.\footnote{The flexibility gained through this in-kernel transpose method removes the necessity of utilizing shuffle instructions to reassign register contents among threads, a requirement discussed previously in~\citep{colfax_fp8_flashattention_2024}.}

\textbf{Accuracy: block quantization and incoherent processing.} In the FP8 (e4m3) representation, only 3 bits are allotted to the mantissa while 4 bits are reserved for the exponent, leading to greater numerical inaccuracies compared to FP16 or BF16 formats. Additionally, large-scale models frequently exhibit outlier values~\citep{dettmers2208llm, sun2024massive} whose magnitudes significantly exceed those of typical entries, which adds challenges to effective quantization. A standard approach involves per-tensor scaling~\citep{micikevicius2022fp8}, where a separate scalar is assigned to each tensor (such as a unique scaling parameter for $\mathbf{Q}$, $\mathbf{K}$, and $\mathbf{V}$). 

To mitigate the increased numerical error found in FP8 computations for attention, we introduce two strategies:

\begin{enumerate}

\item \textbf{Block quantization}: In this approach, a single scalar is maintained for each block, where for each of $\mathbf{Q}$, $\mathbf{K}$, and $\mathbf{V}$, the tensor is partitioned into blocks of dimensions $B_r \times d$ or $B_c \times d$, and each block undergoes separate quantization. This quantization process can be combined with certain operations preceding attention, such as rotary embedding, without incurring any additional slowdown, because rotary embedding is constrained primarily by memory bandwidth. Given that the \textsc{FlashAttention-3} algorithm is inherently block-based, it becomes possible to adjust the scaling of every block in $\mathbf{S}$ to reflect this block-wise quantization without extra computational overhead.

\item \textbf{Incoherent processing}: To mitigate the impact of outliers, both $\mathbf{Q}$ and $\mathbf{K}$ are pre-multiplied by a random orthogonal matrix $\mathbf{M}$ prior to FP8 quantization. The orthogonality of $\mathbf{M}$ ensures that $\mathbf{M} \mathbf{M}^\top = I$, and therefore, $(\mathbf{Q} \mathbf{M}) (\mathbf{K} \mathbf{M})^\top$ is equivalent to $\mathbf{Q} \mathbf{K}^\top$. Thus, applying $\mathbf{M}$ to both $\mathbf{Q}$ and $\mathbf{K}$ leaves the attention computation unchanged. This method effectively disperses outlier values, since each element of $\mathbf{Q} \mathbf{M}$ or $\mathbf{K} \mathbf{M}$ constitutes a random aggregation of the original tensor entries, lowering the quantization error. Following the practices of \citet{chee2024quip} and~\citet{tseng2024quip}, $\mathbf{M}$ is typically formed as the product of random diagonal matrices containing $\pm 1$ and a Hadamard matrix, which enables multiplication in $O(d \log d)$ time instead of $O(d^2)$, and can also be seamlessly integrated with the rotary embedding step, again with no extra computational expense.
\end{enumerate}

Experimental results, presented in \cref{sec:numerical_error}, demonstrate that these two methods can decrease numerical error by as much as a factor of 2.6.

\section{Empirical Validation}
\label{sec:exp}

The implementation and assessment of \textsc{FlashAttention-3} relies on the WGMMA and TMA abstractions provided by CUTLASS~\citep{Thakkar_CUTLASS_2023}, enabling us to measure both its performance and precision.

\begin{itemize}

\item \textbf{Attention benchmarking.} We assess the performance of \textsc{FlashAttention-3} across a range of sequence lengths, comparing its execution time against the default PyTorch implementation, \textsc{FlashAttention-2}, \textsc{FlashAttention-2} running on Triton with H100-optimized instructions, and an H100-tuned \textsc{FlashAttention-2} implementation from cuDNN. Our findings indicate that \textsc{FlashAttention-3} achieves speedups of up to 2.0$\times$ relative to \textsc{FlashAttention-2}, and outpaces the Triton variant by up to 1.5$\times$.
Peak throughput with \textsc{FlashAttention-3} is 740 TFLOPs/s, equivalent to 75\% of the H100 GPU’s theoretical maximum.
\item \textbf{Ablation analysis.} Through systematic ablation, we demonstrate that key improvements—warp-specialization and GEMM-softmax pipeline fusion—directly enable the observed acceleration in \textsc{FlashAttention-3}.
\item \textbf{Evaluation of FP8 attention accuracy.} Experiments establish that the combination of block quantization and incoherent computation lowers numerical errors in FP8 \textsc{FlashAttention-3} by a factor of 2.6$\times$.
\end{itemize}

\subsection{Benchmarking Attention}
\label{subsec:benchmark_attn}

We evaluate the execution time of several attention mechanisms on an H100 80GB SXM5 GPU under a range of configurations (with and without a causal mask, head dimensions set to 64 or 128) using FP16 data types. The benchmark results, presented in~\cref{fig:benchmark_attn_fwd} and~\cref{fig:benchmark_attn_bwd}, indicate that \textsc{FlashAttention-3} achieves a speedup of approximately 1.5-2.0$\times$ over \textsc{FlashAttention-2} in the forward computation, and an improvement of about 1.5-1.75$\times$ in the backward computation. When placed alongside a traditional attention implementation, \textsc{FlashAttention-3} demonstrates acceleration of up to 3-16$\times$. Notably, for sequence lengths of 1k and beyond, \textsc{FlashAttention-3} not only outperforms other software implementations but also exceeds the performance of the proprietary cuDNN library, specifically tuned for the H100 GPU platform.

\paragraph{Benchmark settings:} We vary the sequence length as 512, 1k, ..., 16k, and set batch size so that the total number of tokens is 16k. We set the hidden dimension to 2048, and head dimension to be either 64, 128, or 256 (i.e., 32 heads, 16 heads, or 8 heads). To calculate the FLOPs of the forward pass, we use:

\begin{equation*}
  4 \cdot \text{seqlen}^2 \cdot \text{head dimension} \cdot \text{number of heads}.
\end{equation*}

Applying causal masking entails halving the value to reflect that nearly 50% of the entries undergo computation. For the backward pass FLOPs, the forward pass count is scaled by a factor of 2.5; this adjustment stems from the presence of 2 matrix multiplications in the forward phase and 5 in the backward phase, which is attributed to recomputation.

Additionally, we evaluate the forward pass runtime using FP8 with comparable configurations. The findings for headdim 256 are illustrated in~\cref{fig:benchmark_attn_fp8}, while comprehensive results are presented in \cref{sec:benchmark_attn_fp8_full}.

\subsection{Ablation Study: 2-Stage Pipelining Experiments}
\label{sec:2-stage-pipelining-experiments}

We conduct an ablation study on both the 2-stage WGMMA-softmax pipelining and warp-specialization for non-causal FP16 \textsc{FlashAttention-3}, maintaining fixed parameters $\{\text{batch}, \text{seqlen}, \text{nheads}, \text{hdim}\} = \{ 4, 8448, 16, 128\}$. Results presented in~\cref{table:ablation_pipelining} demonstrate that the algorithmic enhancements—specifically the incorporation of asynchronous execution with warp-specialization and the overlap of GEMM and softmax—produce a marked increase in throughput, rising from 570 to 661 TFLOPs.

\subsection{Numerical Error Validation}
\label{sec:numerical_error}

Given ongoing attention to the numerical inaccuracies associated with \textsc{FlashAttention}~\citep{golden2024flash}, we evaluate the numerical precision of \textsc{FlashAttention-2}, \textsc{FlashAttention-3}, and a conventional attention mechanism by benchmarking each against an FP64 reference implementation. To model the presence of outlier features and activations typical in large language models (LLMs)~\citep{dettmers2208llm, sun2024massive}, we initialize the elements of $\mathbf{Q}, \mathbf{K}, \mathbf{V}$ following the distribution below:

\begin{equation*}
  \mathcal{N}(0, 1) + \mathcal{N}(0, 100) \cdot \mathrm{Bernoulli}(0.001).
\end{equation*}

Specifically, every entry follows a normal distribution with a mean of zero and a standard deviation of 1; however, for 0.1\% of entries, we include an additional, independent Gaussian component with a standard deviation of 10. The root mean squared error (RMSE) is then reported in~\cref{table:numerical_error}.
When using FP16, both \textsc{FlashAttention-2} and \textsc{FlashAttention-3} demonstrate an RMSE that is 1.7$\times$ lower than the standard method, due to retaining intermediate values (such as the softmax output) in FP32 precision. The baseline implementation for attention in FP8 applies a per-tensor scaling strategy, performing matmul accumulation in FP32 and maintaining softmax intermediates in FP16. Owing to the use of block quantization and incoherent processing, \textsc{FlashAttention-3} in FP8 achieves a 2.6$\times$ accuracy improvement relative to the baseline.

\section{Dicussion, Limitations, Conclusion}
\label{sec:discussion}

Through \textsc{FlashAttention-3}, we illustrate how innovative programming paradigms and hardware advancements like asynchrony and reduced-precision computation significantly influence attention’s speed and fidelity. Our results indicate a 1.5-2.0$\times$ improvement in attention performance over \textsc{FlashAttention-2}, as well as a 2.6$\times$ reduction in FP8 numerical error relative to conventional per-tensor quantization methods. However, several aspects remain to be explored in future work: further tuning for LLM inference, the incorporation of a persistent kernel architecture within the FP8 kernel,\footnote{In our benchmarking, FP16 \textsc{FlashAttention-3} benefits from both persistent kernel and a load balancing approach, whereas FP8 \textsc{FlashAttention-3} does not implement these optimizations. This is a partial reason why FP8 \textsc{FlashAttention-3} underperforms relative to FP8 cuDNN kernels for short sequence lengths and causal masking.} and investigating how low-precision attention operates during expansive model training. Although our study centers on Hopper GPUs, we anticipate the underlying strategies will transfer effectively to other compute platforms. Ultimately, we aim for enhanced attention implementations to drive progress on tasks requiring extended contextual reasoning.

\end{document}